\section{Introduction}\label{introduction}

Extracting planar surfaces from three-dimensional point clouds is a fundamental task in computer vision and robotics. Walls, floors, and other approximately planar structures provide a compact description of an environment and geometric constraints for reconstruction, registration, and robot localization. Plane representations are particularly attractive in structured indoor scenes, where large surfaces can replace many individual point measurements in a map \cite{kaess2015planes,schaefer2019planes}. Reliable extraction therefore requires more than an accurate plane equation: the method must also identify the observations that support each surface while rejecting measurements that do not belong to any recovered plane.\footnote{The MATLAB source code, confirmed settings, and reproduction instructions for this comparison are publicly available on GitHub~\cite{jing2026planes}: \url{https://github.com/Mattjing/plane-segmentation}.}

This association is difficult because the supporting points are not known in advance. Measurement noise perturbs surface geometry, occlusion produces incomplete observations, and clutter introduces competing structures and outliers. Near the intersection of two planes, a point may be consistent with more than one model. Unequal surface sizes further complicate estimation because a small plane can be obscured by the support of a larger neighboring surface. Although three non-collinear points determine a plane, a noisy or mixed sample may produce an unreliable hypothesis. Conversely, fitting a model to all observations can combine several distinct surfaces. The problem is a joint optimization over continuous plane parameters and discrete point-to-plane assignments. Unlike a classifier that learns reusable weights, it estimates the models and their supporting point sets anew for each scan.

Random Sample Consensus (RANSAC), introduced by Fischler and Bolles \cite{fischler1981ransac}, addresses robust model estimation by drawing small subsets of observations and evaluating the support of the resulting hypotheses. For plane extraction, three non-collinear sampled points generate a candidate plane. The candidate is tested against all available points using orthogonal point-to-plane distances, and its consensus set contains those within a specified tolerance. Repeated sampling increases the opportunity to obtain a hypothesis from points on the same surface. The candidate with the largest support can then be refined by least-squares fitting to its inliers. This procedure is simple to implement and can tolerate substantial outlier contamination when a sufficiently supported model is present.

Sequential RANSAC extends this procedure to multiple planes by extracting one model, removing its inliers, and repeating the search on the remaining observations. Its direct computation and explicit geometric threshold make it a useful baseline. However, the extraction is greedy: points assigned to an early plane are unavailable to later models, and an erroneous removal can distort the remaining segmentation. Small surfaces may be missed when their support is weak, while threshold choices can cause neighboring surfaces to merge or individual surfaces to fragment. The sampling effort also increases as useful consensus sets become harder to find. These limitations motivate methods that reconsider assignments across several candidate models.

A later alternative is the PEARL framework of Isack and Boykov, published in 2012 \cite{isack2012energy}. PEARL begins with sampled model hypotheses but formulates their estimation as the minimization of a common energy. Its basic formulation combines geometric fitting costs with a spatial regularity term that discourages inconsistent labels among neighboring observations. For fixed models, \(\alpha\)-expansion updates the assignments through graph-cut moves \cite{boykov2001graphcuts}. For fixed assignments, the model parameters are re-estimated from their supporting points. Repeating these steps allows observations to move between competing models and improves the models beyond their original sampled estimates. For the squared orthogonal plane-fitting costs used here, re-estimation is performed by least-squares fitting using singular value decomposition.

The flexibility of joint reassignment introduces additional computational work. In our implementation, each \(\alpha\)-expansion sweep attempts a move for every candidate label, and each move requires constructing and solving a binary cut problem over the point neighborhood graph. These sweeps are repeated as the plane parameters change, alongside the evaluation of fitting costs and the least-squares updates. Runtime consequently depends on the number of points, graph edges, candidate planes, and optimization iterations. The solution also depends on initialization and regularization strength, and global optimality is not guaranteed \cite{isack2012energy}. Spatial coherence can resolve ambiguous assignments, but excessive regularization can suppress valid boundaries. Energy reduction alone therefore does not establish agreement with a reference segmentation.

This paper examines these accuracy and computational tradeoffs through a MATLAB implementation on the Synthetic Plane Extraction Benchmark (SynPEB) \cite{schaefer2019planes}. The study compares sequential RANSAC, PEARL initialized with sampled proposals, and a hybrid that applies PEARL refinement to the planes and assignments produced by RANSAC. It also removes spatial regularization to assess its role under otherwise fixed settings. Hyperparameters are selected using training scans and then held fixed during independent fitting of each test scan; no plane classifier is learned across scans. Evaluation uses reference plane memberships, matched segmentation metrics, repeated random seeds, energy histories, and algorithm runtimes. The aim is to determine when joint optimization improves greedy extraction and how much of the outcome depends on the initial model set.
